\chapter{Revisit 1}

% ================= TOPIC 1 =================
\begin{lessonbox}{1. Simple sentence with です (Kalimat Dasar)}
    Halo Ayah, Bunda, dan Adik-adik! Kata \textbf{です (desu)} adalah kata ajaib yang dipakai di akhir kalimat agar terdengar sopan. Anggap saja artinya "Adalah". Kalau ditanya nama oleh orang Jepang, kamu tinggal sebut namamu dan tambah "desu"!
\end{lessonbox}
\begin{convobox}{1}
    \noindent
    \jpword{わたし}{私}{Watashi} \jpkana{は}{wa} 
    \jpkana{アンディ}{Andi} \jpkana{です。}{desu.}

    \vspace{1.5em}
    \noindent{\Large \color{articolor} \textbf{Arti:} Saya adalah Andi.}
\end{convobox}

% ================= TOPIC 2 =================
\begin{lessonbox}{2. The possessive の (Kepemilikan "Milik")}
    Saat kamu menemukan barangmu yang tertinggal di hotel, kamu bisa bilang "Ini milikku!". Partikel \textbf{の (no)} berfungsi menyambungkan pemilik dan barangnya.
\end{lessonbox}
\begin{convobox}{2}
    \noindent
    \jpword{わたし}{私}{Watashi} \jpkana{の}{no} 
    \jpkana{かばん}{kaban} \jpkana{です。}{desu.}

    \vspace{1.5em}
    \noindent{\Large \color{articolor} \textbf{Arti:} Ini adalah tas milik saya.}
\end{convobox}

% ================= TOPIC 3 =================
\begin{lessonbox}{3. Adjective basics (Kata Sifat Dasar)}
    Wah, makanannya enak! Pemandangannya indah! Untuk memuji sesuatu di Jepang, sebutkan bendanya lalu tambahkan kata sifat dan tutup dengan \textit{desu}.
\end{lessonbox}
\begin{convobox}{3}
    \noindent
    \jpkana{この}{Kono} \jpword{すし}{寿司}{sushi} \jpkana{は}{wa} 
    \jpkana{おいしい}{oishii} \jpkana{です。}{desu.}

    \vspace{1.5em}
    \noindent{\Large \color{articolor} \textbf{Arti:} Sushi ini lezat (enak).}
\end{convobox}

% ================= TOPIC 4 =================
\begin{lessonbox}{4. Adjectives - negative and linking (Kata Sifat: Tidak!)}
    Bagaimana kalau ternyata makanannya TIDAK pedas? Untuk kata sifat berakhiran "i", ubah huruf "i" terakhir menjadi \textbf{〜くない (〜kunai)}. 
\end{lessonbox}
\begin{convobox}{4}
    \noindent
    \jpkana{この}{Kono} \jpkana{カレー}{karee} \jpkana{は}{wa} 
    \jpword{から}{辛}{kara}\jpkana{くない}{kunai} \jpkana{です。}{desu.}

    \vspace{1.5em}
    \noindent{\Large \color{articolor} \textbf{Arti:} Kari ini tidak pedas.}
\end{convobox}

% ================= TOPIC 5 =================
\begin{lessonbox}{5. Adjectives: like, dislike, and want (Suka \& Ingin)}
    Kalau Adik-adik ingin es krim atau mainan baru, kalian bisa gunakan \textbf{好き (Suki - Suka)} atau \textbf{欲しい (Hoshii - Ingin)}. Jangan lupa pakai partikel \textit{ga} ya!
\end{lessonbox}
\begin{convobox}{5}
    \noindent
    \jpword{わたし}{私}{Watashi} \jpkana{は}{wa} 
    \jpkana{アイス}{aisu} \jpkana{が}{ga} 
    \jpword{す}{好}{su}\jpkana{き}{ki} \jpkana{です。}{desu.}

    \vspace{1.5em}
    \noindent{\Large \color{articolor} \textbf{Arti:} Saya suka es krim.}
\end{convobox}

% ================= TOPIC 6 =================
\begin{lessonbox}{6. Simple questions (Bertanya Darurat)}
    Saat jalan-jalan, Ayah dan Bunda pasti butuh bertanya arah. Tambahkan saja \textbf{か (ka)} di akhir kalimat, dan kalimatmu otomatis berubah menjadi pertanyaan!
\end{lessonbox}
\begin{convobox}{6}
    \noindent
    \jpkana{トイレ}{Toire} \jpkana{は}{wa} 
    \jpkana{どこ}{doko} \jpkana{ですか？}{desu ka?}

    \vspace{1.5em}
    \noindent{\Large \color{articolor} \textbf{Arti:} Toilet ada di mana?}
\end{convobox}

% ================= TOPIC 7 =================
\begin{lessonbox}{7. Basic Verbs: ます and ません (Kata Kerja Dasar)}
    Untuk bercerita tentang apa yang akan kamu lakukan, gunakan akhiran \textbf{〜ます (〜masu)}. Kalau TIDAK akan melakukannya, ubah menjadi \textbf{〜ません (〜masen)}.
\end{lessonbox}
\begin{convobox}{7}
    \noindent
    \jpkana{ラーメン}{Raamen} \jpkana{を}{o} 
    \jpword{た}{食}{ta}\jpkana{べます。}{bemasu.}

    \vspace{1.5em}
    \noindent{\Large \color{articolor} \textbf{Arti:} Saya (akan) makan ramen.}
\end{convobox}

% ================= TOPIC 8 =================
\begin{lessonbox}{8. Basic particles (を dan で)}
    Partikel adalah lem yang merekatkan kalimat! \textbf{を (o)} menempel pada objek yang kamu makan/beli. Sedangkan \textbf{で (de)} menempel pada TEMPAT kamu melakukan aktivitas itu.
\end{lessonbox}
\begin{convobox}{8}
    \noindent
    \jpkana{ホテル}{Hoteru} \jpkana{で}{de} 
    \jpword{あさ}{朝}{asa}\jpkana{ごはんを}{gohan o} 
    \jpword{た}{食}{ta}\jpkana{べます。}{bemasu.}

    \vspace{1.5em}
    \noindent{\Large \color{articolor} \textbf{Arti:} Saya makan sarapan di hotel.}
\end{convobox}

% ================= TOPIC 9 =================
\begin{lessonbox}{9. Basic particles (へ/に - Arah \& Tujuan)}
    Saat naik kereta, perhatikan partikel \textbf{へ (e)} atau \textbf{に (ni)}. Keduanya berarti "Ke". Ini sangat penting untuk memberi tahu supir taksi ke mana keluargamu mau pergi!
\end{lessonbox}
\begin{convobox}{9}
    \noindent
    \jpword{きょうと}{京都}{Kyouto} \jpkana{に}{ni} 
    \jpword{い}{行}{i}\jpkana{きます。}{kimasu.}

    \vspace{1.5em}
    \noindent{\Large \color{articolor} \textbf{Arti:} Saya pergi ke Kyoto.}
\end{convobox}

% ================= TOPIC 10 =================
\begin{lessonbox}{10. Basic particles (も/と - Juga \& Dan)}
    Saat berbelanja suvenir, kadang kita mau beli banyak hal. Gunakan \textbf{と (to)} untuk bilang "Dan", serta \textbf{も (mo)} untuk bilang "Juga".
\end{lessonbox}
\begin{convobox}{10}
    \noindent
    \jpkana{これ}{Kore} \jpkana{と}{to} 
    \jpkana{それ}{sore} \jpkana{を}{o} 
    \jpword{くだ}{下}{kuda}\jpkana{さい。}{sai.}

    \vspace{1.5em}
    \noindent{\Large \color{articolor} \textbf{Arti:} Tolong (berikan saya) ini dan itu.}
\end{convobox}

% ================= TOPIC 11 =================
\begin{lessonbox}{11. Existence Verbs: ある/いる (Ada!)}
    Di taman ada rusa? Di toko ada payung? Gunakan \textbf{います (Imasu)} jika yang ada adalah makhluk hidup (manusia/hewan). Gunakan \textbf{あります (Arimasu)} untuk benda mati (gedung/barang).
\end{lessonbox}
\begin{convobox}{11}
    \noindent
    \jpkana{あそこ}{Asoko} \jpkana{に}{ni} 
    \jpkana{パンダ}{panda} \jpkana{が}{ga} 
    \jpkana{います。}{imasu.}

    \vspace{1.5em}
    \noindent{\Large \color{articolor} \textbf{Arti:} Di sana ada Panda.}
\end{convobox}

% ================= TOPIC 12 =================
\begin{lessonbox}{12. Past tense: でした (Masa Lalu: Benda/Sifat)}
    Liburan kemarin hujan? Atau makanannya sangat enak? Jika peristiwanya SUDAH BERLALU, kata \textit{desu} diubah menjadi \textbf{でした (deshita)}.
\end{lessonbox}
\begin{convobox}{12}
    \noindent
    \jpword{きのう}{昨日}{Kinou} \jpkana{は}{wa} 
    \jpword{あめ}{雨}{ame} \jpkana{でした。}{deshita.}

    \vspace{1.5em}
    \noindent{\Large \color{articolor} \textbf{Arti:} Kemarin (cuacanya) hujan.}
\end{convobox}

% ================= TOPIC 13 =================
\begin{lessonbox}{13. Past tense: ました (Masa Lalu: Kata Kerja)}
    Sama seperti poin sebelumnya, kalau kegiatan (kata kerja) sudah selesai kamu lakukan hari itu, ubah akhiran \textit{~masu} menjadi \textbf{〜ました (〜mashita)}.
\end{lessonbox}
\begin{convobox}{13}
    \noindent
    \jpkana{たこやき}{Takoyaki} \jpkana{を}{o} 
    \jpword{か}{買}{ka}\jpkana{いました。}{imashita.}

    \vspace{1.5em}
    \noindent{\Large \color{articolor} \textbf{Arti:} Saya sudah membeli Takoyaki.}
\end{convobox}

% ================= TOPIC 14 =================
\begin{lessonbox}{14. Invitations (Mari Mengajak!)}
    Ayah dan Bunda mau mengajak anak-anak pergi ke Disneyland? Ubah kata kerjanya menjadi \textbf{〜ましょう (〜mashou)} yang artinya "Mari kita... / Ayo kita...".
\end{lessonbox}
\begin{convobox}{14}
    \noindent
    \jpword{いっしょ}{一緒}{Issho} \jpkana{に}{ni} 
    \jpword{い}{行}{i}\jpkana{きませんか？}{kimasen ka?}

    \vspace{1.5em}
    \noindent{\Large \color{articolor} \textbf{Arti:} Maukah pergi bersama-sama?}
\end{convobox}

% ================= TOPIC 15 =================
\begin{lessonbox}{15. Verb form: て (Bentuk Te untuk Menyambung)}
    Bentuk \textit{Te} sangat ajaib! Bisa dipakai untuk menyambung kalimat ("Saya makan, \textit{lalu} pergi") atau meminta tolong secara kasual.
\end{lessonbox}
\begin{convobox}{15}
    \noindent
    \jpkana{ちょっと}{Chotto} \jpword{ま}{待}{ma}\jpkana{って}{tte} 
    \jpword{くだ}{下}{kuda}\jpkana{さい。}{sai.}

    \vspace{1.5em}
    \noindent{\Large \color{articolor} \textbf{Arti:} Tolong tunggu sebentar.}
\end{convobox}

% ================= TOPIC 16 =================
\begin{lessonbox}{16. Verb form: ている (Sedang Berlangsung)}
    Jika saat ini kamu SEDANG melakukan sesuatu (misalnya sedang duduk di kereta), gunakan bentuk \textbf{〜ています (〜te imasu)}.
\end{lessonbox}
\begin{convobox}{16}
    \noindent
    \jpword{いま}{今}{Ima}\jpkana{、}{,} 
    \jpword{でんしゃ}{電車}{densha} \jpkana{に}{ni} 
    \jpword{の}{乗}{no}\jpkana{っています。}{tte imasu.}

    \vspace{1.5em}
    \noindent{\Large \color{articolor} \textbf{Arti:} Sekarang, saya sedang naik kereta.}
\end{convobox}

% ================= TOPIC 17 =================
\begin{lessonbox}{17. Requests and permission (Meminta Izin Terdengar Sopan)}
    Saat melihat maskot lucu atau pemandangan indah di kuil, kadang kita tidak yakin boleh memfotonya atau tidak. Minta izinlah dengan sopan pakai \textbf{〜てもいいですか (〜te mo ii desu ka?)}.
\end{lessonbox}
\begin{convobox}{17}
    \noindent
    \jpword{しゃしん}{写真}{Shashin} \jpkana{を}{o} 
    \jpword{と}{撮}{to}\jpkana{ってもいいですか？}{tte mo ii desu ka?}

    \vspace{1.5em}
    \noindent{\Large \color{articolor} \textbf{Arti:} Bolehkah saya mengambil foto?}
\end{convobox}

% --- SUMMARY BOX CLOSING ---
\vspace{2em}
\begin{tcolorbox}[
    colback=green!5!white,
    colframe=green!60!black,
    boxrule=3pt,
    arc=5mm,
    title={\huge\bfseries \sffamily 🌸 Selesai! Selamat Berlibur! 🌸},
    fonttitle=\color{white},
    attach boxed title to top center={yshift=-4mm},
    boxed title style={colback=green!60!black, rounded corners},
    left=5mm, right=5mm, top=7mm, bottom=5mm
]
\Large Keluarga Hebat! Kalian telah menyelesaikan seluruh pelajaran penting dalam buku \textbf{Andi - Japanese Handbook} ini. Dengan bekal tata bahasa dasar dan kumpulan kosakata penting ini, pengalaman berlibur kalian di Jepang dijamin akan menjadi jauh lebih seru, aman, dan berkesan. Jangan takut mencoba berbicara dengan penduduk lokal ya. \textbf{日本の旅行を楽しんでください！ (Nihon no ryokou o tanoshinde kudasai - Selamat menikmati liburan di Jepang!)}
\end{tcolorbox}

\newpage